\section*{Capítulo \ref{ch:g2:from-seed-to-plant} --- De la semilla a la planta}

\begin{solution}{exo:g2:from-seed-to-plant:1}
Una plantita ya empezada (el germen de la \omterm{def:g1:plants-around-us:plant}{planta}, con el principio de
una \omterm{def:g1:plants-around-us:parts}{raíz} y de un brote) y una despensa de alimento para mantenerla.
\end{solution}

\begin{solution}{exo:g2:from-seed-to-plant:2}
\omterm{def:g2:from-seed-to-plant:germination}{Germinar} es despertarse y empezar a crecer: la \omterm{def:g2:from-seed-to-plant:seed}{semilla} se hincha, su
cubierta se abre y empieza a convertirse en \omterm{def:g1:plants-around-us:plant}{planta}. La \omterm{def:g1:plants-around-us:parts}{raíz} sale la
primera y se hunde hacia abajo.
\end{solution}

\begin{solution}{exo:g2:from-seed-to-plant:3}
Agua y calor. La luz no está en la lista --- la \omterm{def:g2:from-seed-to-plant:seed}{semilla} \omterm{def:g2:from-seed-to-plant:germination}{germina} tan
tranquila a oscuras, bajo tierra.
\end{solution}

\begin{solution}{exo:g2:from-seed-to-plant:4}
Vive de la despensa que lleva dentro la \omterm{def:g2:from-seed-to-plant:seed}{semilla} --- las dos mitades
gordas de la judía. Esa despensa alimenta a la \omterm{prop:g2:from-seed-to-plant:cycle}{plántula} hasta que sus
\omterm{def:g1:plants-around-us:parts}{hojas} verdes llegan a la luz y puede alimentarse sola.
\end{solution}

\begin{solution}{exo:g2:from-seed-to-plant:5}
Por ejemplo: judías, guisantes, lentejas, arroz, pipas de girasol,
pepitas de manzana, un hueso de melocotón.
\end{solution}

\begin{solution}{exo:g2:from-seed-to-plant:6}
Porque una \omterm{def:g2:from-seed-to-plant:seed}{semilla} necesita calor además de agua: en la tierra fría
del invierno esperaría, o se pudriría, en vez de \omterm{def:g2:from-seed-to-plant:germination}{germinar}. La
primavera trae la tierra templada que la \omterm{def:g2:from-seed-to-plant:seed}{semilla} está esperando.
\end{solution}

\begin{solution}{exo:g2:from-seed-to-plant:7}
\omterm{def:g2:from-seed-to-plant:seed}{Semilla}, \omterm{prop:g2:from-seed-to-plant:cycle}{plántula}, mata frondosa, \omterm{def:g1:plants-around-us:plant}{planta} en \omterm{def:g1:plants-around-us:parts}{flor}, vaina con judías nuevas.
\end{solution}

\begin{solution}{exo:g2:from-seed-to-plant:8}
Las dos \omterm{def:g1:plants-around-us:parts}{raíces} crecen hacia abajo y los dos brotes hacia arriba, sin
importar cómo estuvieran colocadas las judías. La \omterm{prop:g2:from-seed-to-plant:cycle}{plántula} siempre
encuentra el arriba y el abajo --- justo lo que predice la observación.
\end{solution}

\begin{solution}{exo:g2:from-seed-to-plant:9}
En el cajón las \omterm{def:g2:from-seed-to-plant:seed}{semillas} no tenían agua, así que no podían \omterm{def:g2:from-seed-to-plant:germination}{germinar};
pero una \omterm{def:g2:from-seed-to-plant:seed}{semilla} buena sin agua no se muere --- espera. Un año después,
con agua y calor, las \omterm{def:g2:from-seed-to-plant:seed}{semillas} que esperaban se despertaron con normalidad.
\end{solution}

\begin{solution}{exo:g2:from-seed-to-plant:10}
La despensa de la \omterm{def:g2:from-seed-to-plant:seed}{semilla} --- las dos mitades gordas de la judía ---
hace el papel de la yema. El \omterm{def:g2:animals-grow-up:born}{huevo} y la \omterm{def:g2:from-seed-to-plant:seed}{semilla} son los dos comienzos
de un \omterm{def:g1:living-or-not:living}{ser vivo} nuevo, con el alimento que la cría necesita antes de
poder alimentarse sola.
\end{solution}
